\clearpage\thispagestyle{empty}
\pdfbookmark[0]{Quellen zu Gauss an Bode (21. März 1820)}{bodemarch1820-sources}
\begin{center}\textbf{Quellen zu Gauss an Bode (21. März 1820)}\end{center}
\noindent Diese KI-Übertragung ins heutige Deutsch folgt \textit{Gauss an Bode} in den \href{https://gdz.sub.uni-goettingen.de/id/PPN236007467}{\textit{Werken}, Band~VI} (Göttingen, 1874), S.~428--429. Die \href{https://api.digitale-sammlungen.de/iiif/image/v2/bsb10538352_00166/full/1600,/0/default.jpg}{Überschrift des Erstdrucks} im von Johann Elert Bode herausgegebenen und 1820 in Berlin erschienenen \textit{Astronomischen Jahrbuch für das Jahr 1823} nennt Beobachtungen von Saturn und Vesta im Jahr 1819 sowie von Pallas und Mars im Jahr 1820. Sie wurden am Mittagsfernrohr der Göttinger Sternwarte gemacht und von Gauß am 21. März 1820 eingesandt. Im Exemplar der Bayerischen Staatsbibliothek steht der Bericht auf S.~160--161.
\par\smallskip\noindent Die Werkausgabe gibt als letzten Zahlenwert für den Mars am 23. Januar 8.0 an, der \href{https://api.digitale-sammlungen.de/iiif/image/v2/bsb10538352_00167/full/1600,/0/default.jpg}{Erstdruck} 8,1. Im \hyperlink{native-427-part-lower}{Instrumentenbericht} vom Oktober 1819 steht für Saturn am 1. Oktober 36.3 und für Vesta am 24. September 58.0. Beide Ausgaben des Märzbriefs geben dafür 36.9 bzw. 58.1 an. Die Tabellen folgen auch in der Anordnung der Einheiten und Überschriften der Werkausgabe. KI-Übertragung und KI-Satz.
\clearpage
